\begin{qfdefinition}[Real estate: features and characteristics of the asset, forms of ownership, and its investment characteristics]\label{def:real-estate-characteristics}
\qffield{Statement} \emph{Features.} A property is land plus a building: heterogeneous, fixed in location, indivisible and durable. Its value comes from the cash flows it can earn, chiefly rents, net of operating costs. The building depreciates and has a finite economic life, but the land does not \citep[pp.~739--740]{investment-valuation-2012}. The market has three segments: occupiers who lease space, developers who supply it, and investors who price it, with demand measured by take-up and net absorption and supply by stock, vacancy and construction \citep[pp.~5--15]{tsolacos-real-estate-2021}.

\emph{Forms of ownership.}
\begin{itemize}
\item \emph{direct private equity}: owning buildings outright or through private funds;
\item \emph{private debt}: mortgages and mezzanine loans;
\item \emph{public equity}: real estate investment trusts (REITs) and listed property companies; US REITs must distribute at least 90 percent of taxable income \citep[p.~376]{s7-ang};
\item \emph{public debt}: commercial mortgage-backed securities.
\end{itemize}

\emph{Investment characteristics.}
\begin{itemize}
\item \emph{Income}: a large, stable share of return. In a quarterly NCREIF sample the income return is steady while capital growth drives volatility \citep[pp.~17--18]{tsolacos-real-estate-2021}. In the second quarter of 2026 the NCREIF Property Index returned 1.29 percent, 1.17 points of it income \citep[p.~1]{ncreif-npi-2026q2}.
\item \emph{Illiquidity and large lot sizes}, high transaction costs, and dependence on local knowledge \citep[pp.~741, 747]{investment-valuation-2012}.
\item \emph{Appraisal-based returns}: direct real estate returns are smoothed and biased by selection. The NCREIF series has a quarterly autocorrelation of 78 percent and volatility of 2.2 percent, against about 8--9 percent for REITs and equities (1978--2011), so reported diversification overstates the true \citep[pp.~376--378]{s7-ang}.
\item \emph{Partial inflation protection}: residential real estate was the only asset class Fama and Schwert found positively affected by unexpected inflation (Damodaran, p.~740). Ang finds real estate ``only partially real'': correlations with inflation of about 0.3--0.4 at long horizons and little protection at short ones (pp.~380--381).
\item \emph{Leverage} is common in direct and REIT ownership.
\end{itemize}

\qffield{Practice} \emph{Marlow House} (original; \hyperref[meth:real-estate-valuation-indexes]{valuation}): 200,000 square feet of offices let at 40 a square foot at 93 percent occupancy earns NOI of 4.83 million after taxes, operating costs and a capital reserve. A buyer at a 6 percent capitalisation rate pays 80.55 million, about 6 percent income before growth.

\qffield{Conditions} Characteristics described are for income-producing commercial property; owner-occupied housing and development land behave differently.

\qffield{Pitfalls} Treating NCREIF volatility as the risk of owning property. Assuming real estate hedges inflation over short horizons.

\qffield{See also} \hyperref[meth:real-estate-value-drivers]{value drivers}; \hyperref[meth:reit-valuation]{REITs}; \hyperref[meth:private-capital-diversification]{smoothing}.

\end{qfdefinition}

\begin{qfmethod}[Real estate features for valuation, economic value drivers and portfolio role, and the distinctive characteristics of commercial property types]\label{meth:real-estate-value-drivers}
\qffield{Statement} \emph{Features that matter for valuation} \citep[pp.~750--752, 759--761]{investment-valuation-2012}: location, size and age; the lease terms (length, rent reviews, expense reimbursements and stops, renewal options); occupancy and tenant credit; property taxes and operating costs; the land's own value and the building's remaining life; and whether the property is income-producing at all. Comparables must be adjusted for differences in income, size, scale, location, age and construction quality.

\emph{Economic value drivers.} Value is NOI divided by a capitalisation rate (\hyperref[meth:real-estate-valuation-indexes]{valuation}), so value moves with:
\begin{itemize}
\item \emph{rents and occupancy}, driven by occupier demand (employment, consumption, trade) against new supply: take-up, net absorption, vacancy and construction \citep[pp.~5--15]{tsolacos-real-estate-2021};
\item \emph{expected rent growth}, close to inflation in a stable market and above it while vacancy is low \citep[p.~751]{investment-valuation-2012};
\item \emph{the required return}, which follows real interest rates, the risk premium and credit availability; Damodaran treats the cap rate as a required return less growth (p.~752).
\end{itemize}
In a portfolio, real estate gives a large income component, some long-run inflation protection and, through REITs, an equity-like exposure; direct holdings diversify less than their appraisals suggest \citep[pp.~376--381]{s7-ang}.

\emph{Property types} (Interpretation, from Lynn's 2008 sector reviews and NCREIF weights):
\begin{itemize}
\item \emph{office}: volatile, with demand tied to job growth and financial-sector hiring \citep[p.~13]{lynn-real-estate-2009};
\item \emph{retail}: long leases and anchor tenants, exposed to consumer spending and structural change in shopping (pp.~14--15);
\item \emph{industrial and logistics}: easy to build, so supply responds quickly; demand follows trade and distribution; the largest NCREIF sector, 33.3 percent of index value in 2026 \citep[p.~1]{ncreif-npi-2026q2};
\item \emph{multifamily}: demand from demographics and household formation (Lynn, p.~11), with short leases that reprice rents quickly (Interpretation);
\item \emph{hotels}: effectively nightly leases, so revenue per available room tracks the cycle immediately; the most operating-intensive type (pp.~16--17).
\end{itemize}

\qffield{Practice} \emph{Marlow House} (original). Raising occupancy from 93 to 95 percent adds 0.14 million to NOI, 2.4 million of value at a 6 percent cap rate. A 50 basis point rise in the cap rate, from tighter credit or higher real rates, removes 6.2 million: the required return is the larger driver over a cycle.

\qffield{Conditions} Lynn's sector assessments date from 2008 and describe the US; sector risk has since changed, for example office demand after remote working.

\qffield{Pitfalls} Projecting rent growth above inflation forever. Ignoring that short-lease types reprice in both directions.

\qffield{See also} \hyperref[def:real-estate-characteristics]{characteristics}; \hyperref[meth:private-real-estate-underwriting]{private underwriting}.

\end{qfmethod}

\begin{qfmethod}[Due diligence and valuation of real estate: income, sales comparison and cost approaches, and the construction and biases of real estate indexes]\label{meth:real-estate-valuation-indexes}
\qffield{Statement} \emph{Due diligence} verifies what the cash flows rest on: the leases and tenant credit, vacancy and credit loss, operating expenses, reimbursements and expense stops, and property taxes \citep[pp.~750--751]{investment-valuation-2012}; and, in practice, the physical condition and capital needs, title, zoning and environmental issues (Interpretation).

\emph{Approaches.}
\begin{itemize}
\item \emph{Income: direct capitalisation.} $V=\mathrm{NOI}_1/\text{cap}$. With constant growth $g$ and required return $r$, the cap rate on next year's NOI is $r-g$, and on current NOI $(r-g)/(1+g)$ (p.~752).
\item \emph{Income: discounted cash flow.} Projected rents, vacancy, credit loss, expenses and capital spending over a holding period, plus a terminal value, discounted at a cost of capital for property-level flows or a cost of equity for flows after debt (pp.~749--752); discount rates can come from a risk-and-return model or from investor surveys, which Damodaran finds less reliable (pp.~748--749).
\item \emph{Sales comparison.} Prices per square foot, gross income multipliers or price-to-NOI of comparable sales, adjusted for differences; comparables may work better for property than for stocks, because nearby buildings share growth and risk (pp.~759--761).
\item \emph{Cost.} Land value plus the cost of a modern equivalent building, less depreciation for physical deterioration, functional obsolescence (outdated design or excess operating cost against the modern equivalent) and economic or external obsolescence (weak demand, oversupply, an occupier unable to pay market rent) \citep[paras~80.2--80.5, 90.2--90.7]{ivsc-ivs105-exposure-draft-2016}. It rests on substitution: a buyer pays no more than the cost of an asset of equal utility, unless time, risk or other factors intervene (para.~70.1). It suits assets that could be recreated quickly and do not earn income, and otherwise needs corroboration (paras~70.2--70.3).
\end{itemize}

\emph{Indexes.}
\begin{itemize}
\item \emph{Appraisal-based} indexes such as the NCREIF Property Index use quarterly property returns weighted by market value and before fees, from appraisals of properties held by tax-exempt fiduciaries. Income return is NOI over beginning value plus half of capital improvements, less half of partial sales and a third of NOI; appreciation uses the same denominator \citep[p.~3]{ncreif-npi-2026q2}.
\item \emph{Transaction-based} indexes use sales, often by repeat-sales methods \citep[p.~18]{tsolacos-real-estate-2021}.
\item \emph{Biases}: appraisal smoothing and lag, selection (only some properties sell), unlevered appraisal indexes against levered REIT indexes, and indexes nobody can buy \citep[pp.~377--378, 427--428]{s7-ang}.
\end{itemize}

\qffield{Practice} \emph{Marlow House} (original, \nolinkurl{trading_notes.numerical.real_asset_investor.real_assets_checks}):
\begin{center}\small
\begin{tabular}{@{}lr@{}}
\toprule
Approach & Value (EUR m) \\
\midrule
Direct capitalisation, NOI 4.83 at 6.0\% & 80.55 \\
DCF at 8\%, ten years, exit cap 6.5\% less 2\% costs & 81.15 \\
Sales comparison, 403 a sq ft (occupancy-adjusted) & 79.57 \\
Cost: land 25 + 330 a sq ft less 20\% depreciation & 77.80 \\
\bottomrule
\end{tabular}
\end{center}
The exit value is 55 percent of the DCF. With constant 2.5 percent growth and an exit cap of $8\%-2.5\%$, the DCF equals NOI over 5.5 percent exactly, which shows that the cap rate embeds growth. Reconciled value: about 80--81 million. The cost approach agrees here because Marlow's market supports building it: its income value exceeds its cost, a Tobin's Q above one.

\emph{Harbour Point} (original), an older office in an oversupplied submarket. Land 10 plus a modern equivalent at 60, less 20 percent physical deterioration, gives 58. NOI of 3.0 at a 7.5 percent cap gives 40, and comparable sales 42; reconciled value 41. The 17 gap is obsolescence that physical depreciation alone misses: 0.4 a year of excess operating cost capitalised at 7.5 percent (5.33, functional) and 11.67 of economic obsolescence. Depreciated replacement cost is therefore not a floor: market value equals replacement cost only when the asset earns its required return, and falls below it when it earns less \citep[pp.~537--538]{investment-valuation-2012}. Nor is cost a ceiling: when building takes years or permits are scarce, buyers pay above cost for a building that exists (IVS 105, paras~70.1, 70.3).

\emph{Index arithmetic.} A property with beginning value 100, NOI 1.2, capital improvements 0.5 and ending value 101 has an NCREIF income return of 1.2018 percent and appreciation of 0.5008 percent. Appraisals that blend half of the new evidence with the prior value peak a period late (112.5 against a transaction peak of 120) and bottom at 100.0 against 95.

\qffield{Conditions} Damodaran's examples are US office buildings valued in 2000. The cost-approach paragraphs are from the IVSC's 2016 exposure draft of IVS 105, which became part of IVS 2017; later editions renumber the chapter. Harbour Point's split of obsolescence is illustrative. In his Illustration 26.4 the regression intercept is printed as $-2{,}535.50$ in the equation and $-2{,}535.70$ in its use; the printed estimate of 179.46 a square foot follows from $-2{,}535.50$ (p.~761).

\qffield{Pitfalls} Applying a market cap rate to in-place NOI with below-market leases. Treating depreciated replacement cost as a floor (or a ceiling) on value without measuring obsolescence. Treating appraisal indexes as transaction prices. Double counting growth in a cap rate and in projected NOI.

\qffield{See also} \hyperref[meth:real-estate-value-drivers]{value drivers}; \hyperref[meth:reit-valuation]{REIT NAV}.

\end{qfmethod}

\begin{qfmethod}[Publicly traded real estate: types of securities, NAV per share from forward cash NOI, FFO and AFFO, REIT valuation by NAV, price multiples and discounted cash flow, and listed against private real estate]\label{meth:reit-valuation}
\qffield{Statement} \emph{Types.} Equity REITs owning properties, mortgage REITs holding real estate debt, listed property companies that retain earnings and pay tax, and listed real estate funds \citep[pp.~376, 380]{s7-ang}.

\emph{NAV per share.} Capitalise forward cash NOI (excluding straight-line rent and other non-cash items) at a market cap rate, add other assets (cash, development at cost, land), subtract debt and other liabilities, and divide by shares. NAV approximates what the assets would fetch privately, so a share price far from NAV signals either mispricing or a disagreement between public and private markets (Interpretation, using the cap-rate valuation above).

\emph{FFO and AFFO.} Nareit defines \emph{funds from operations} as GAAP net income excluding depreciation and amortisation of real estate, gains and losses on sales of certain real estate, gains and losses from change of control and certain impairments of depreciable real estate \citep[p.~2]{nareit-ffo-2018}. FFO adjusts for historical-cost depreciation, which overstates the economic decline of property. \emph{Adjusted} FFO, cash available for distribution and similar measures also deduct items such as straight-line rent and recurring capital spending. Nareit sees no consensus on a single definition and asks issuers to reconcile any such measure to FFO and GAAP (p.~7). SEC staff accept FFO per share; other adjustments may bar per-share presentation (pp.~3, 7).

\emph{Valuation methods.} NAVPS; relative value by price to FFO or AFFO against peers; and DCF of dividends or AFFO.

\emph{Listed against private} \citep[pp.~376--380]{s7-ang}: REITs offer liquidity, small lot sizes, professional management and transparent prices, but they are levered, more volatile and more correlated with equities. Their returns diverge from direct property in the short run but share long-horizon components, and public markets lead private ones in price discovery. Legal constraints (90 percent payout, income tests) and management fees create agency costs.

\qffield{Practice} \emph{Kestrel REIT} (original):
\begin{itemize}
\item \emph{NAV}: forward cash NOI 80 at a 6.25 percent cap is 1,280; plus cash 40 and development at cost 60, less debt 500 and other liabilities 30, gives 850, or 17.00 a share. At 15.30 the shares trade 10 percent below NAV.
\item \emph{FFO}: net income 45, plus real estate depreciation 50, less a gain on sale of 8, is 87 (1.74 a share). \emph{AFFO}: less straight-line rent of 4 and recurring capital and leasing spending of 15, is 68 (1.36).
\item \emph{Relative value}: 10 times FFO gives 17.40, 12.5 times AFFO 17.00.
\item \emph{DCF}: dividends of 80 percent of AFFO growing 3 percent for five years and 2 percent thereafter, at 9 percent, give 16.55.
\end{itemize}
\emph{Conclusion.} Three methods put value at 16.55--17.40 against a price of 15.30. The discount is worth buying only if the private cap rate of 6.25 percent is not about to rise; at 6.75 percent NAV falls to 15.10 a share, below the price.

\qffield{Conditions} The FFO definition is Nareit's 2018 restatement, effective December 2018; AFFO is company-defined.

\qffield{Pitfalls} Comparing P/FFO across REITs with different capital needs. Using in-place NOI with straight-line rents for NAV. Treating a NAV discount as a margin of safety when private values lag public ones.

\qffield{See also} \hyperref[meth:real-estate-valuation-indexes]{property valuation}; \hyperref[meth:earnings-management-presentation]{non-GAAP measures}.

\end{qfmethod}

\begin{qfmethod}[Underwriting private real estate: core, value-added and opportunistic strategies, value drivers, due diligence and a levered valuation with sensitivity]\label{meth:private-real-estate-underwriting}
\qffield{Statement} \emph{Private real estate features.} Private funds buy, transform and sell properties; strategies run from core to opportunistic as transformation and leverage rise \citep[pp.~137--141]{demaria-private-markets-2021}:
\begin{center}\small
\begin{tabular}{@{}lrrr@{}}
\toprule
Strategy & Leverage (\% of value) & Income share of return & Pooled net IRR, US \\
\midrule
Core & 15--30 & 90--100\% & 6.0\% (core and core-plus) \\
Core-plus & 30--50 & 80--90\% & \\
Value-added & 40--70 & 30--50\% & 6.1\% \\
Opportunistic & 60--80 & 0--10\% & 7.9\% \\
\bottomrule
\end{tabular}
\end{center}
Value-added and opportunistic performance is dispersed (bottom-5-percent value-added funds $-10.3$ percent, top 28.9 percent), and the authors expect 8--10 and 12--18 percent respectively over the long term (pp.~140--141).

\emph{Value drivers} are those of property (\hyperref[meth:real-estate-value-drivers]{value drivers}) plus the manager's plan: leasing up, repositioning, refurbishment, cost control and the exit cap rate. \emph{Due diligence} adds the business plan's capital budget and timing, the loan terms and the manager's record.

\emph{Valuation.} Unlevered value from property cash flows at the property's required return; levered returns to equity after debt service, judged against the risk the leverage adds (Interpretation).

\qffield{Practice} \emph{Ferris Park} (original, value-added industrial). Buy for 50 with NOI of 2.5 and spend 6 at once; NOI rises to 4.0 in year 3 and 3 percent a year after; sell after year 5 at a 5.75 percent cap on year-6 NOI, less 1.5 percent costs; 30 of interest-only debt at 6 percent:
\begin{center}\small
\begin{tabular}{@{}lrr@{}}
\toprule
 & Exit cap 5.75\% & Exit cap 6.25\% \\
\midrule
Exit value & 74.88 & 68.89 \\
Unlevered IRR & 11.69\% & 10.13\% \\
Levered IRR (equity 26) & 16.96\% & 14.16\% \\
\bottomrule
\end{tabular}
\end{center}
Interest cover is 1.39 times in year 1 and 2.22 times by year 3; loan to value is 60 percent at purchase. \emph{Decision}: the deal fits a value-added target of 8--10 percent unlevered only if the lease-up happens. The levered return falls 2.8 points for a 50 basis point higher exit cap, so the bid should not exceed 50 unless the exit assumption is conservative.

\qffield{Conditions} The return statistics are Cambridge Associates and eFront data quoted by Demaria and co-authors for vintages to 2008, affected by the 2007--2009 downturn.

\qffield{Pitfalls} Underwriting value-added returns on core exit cap rates. Letting leverage, not the business plan, produce the target return.

\qffield{See also} \hyperref[meth:real-assets-allocation]{allocation}; \hyperref[meth:private-debt-credit-analysis]{lender's view}.

\end{qfmethod}

\begin{qfmethod}[Private real estate, timberland and farmland in a strategic allocation: risk and return against other investments, and a recommendation]\label{meth:real-assets-allocation}
\qffield{Statement} \emph{Private real estate} offers income, partial long-run inflation protection and diversification that appraisals overstate; adjusted for leverage, sector mix and smoothing, direct and listed returns may not differ much \citep[pp.~376--381]{s7-ang}. Fund returns fall with transformation risk less than leverage suggests (\hyperref[meth:private-real-estate-underwriting]{underwriting}).

\emph{Timberland and farmland} give yield from harvests or rents and appreciation from land and, for timber, biological growth. Their income depends on commodity prices rather than leases, and they offer some inflation hedge \citep[p.~143]{demaria-private-markets-2021}. In 2022 the NCREIF Timberland Index returned 12.90 percent; in the first quarter of 2023 it returned 1.75 percent while the NCREIF Property Index lost 1.81 percent \citep[p.~28]{ncreif-timberland-2023q1}. Damodaran's table pairs NCREIF farmland with commercial property and REITs as appraisal-based series \citep[p.~743, Table~26.1]{investment-valuation-2012}.

\emph{Allocation} (Interpretation): decide the real-asset sleeve by the liabilities it supports (inflation-linked or long-dated), the liquidity it can tolerate and the access available; within it, core property and core infrastructure for income, value-added for return, and land for inflation and diversification from equity-driven property.

\qffield{Practice} \emph{A pension with inflation-linked liabilities} (original): a 12 percent real-asset sleeve of 6 percent core and core-plus property (open-ended funds and REITs for liquidity), 2 percent value-added funds, 3 percent core infrastructure (\hyperref[meth:infrastructure-allocation]{infrastructure}) and 1 percent timberland and farmland through a specialist manager. Risk is set on unsmoothed volatility, with REITs used to rebalance quickly. Farmland at the 2026 US average of 6,020 dollars an acre and 160 of cash rent yields 2.66 percent gross, so its return depends on rent and land growth (\hyperref[meth:land-timber-farmland]{land}).

\qffield{Conditions} Timberland figures are from an NCREIF member publication for a single quarter and year; they illustrate, not establish, low correlation.

\qffield{Pitfalls} Counting appraisal-based diversification twice, once in property and again in land. Ignoring that farmland and timber yields are low and their returns depend on appreciation.

\qffield{See also} \hyperref[meth:alternatives-allocation-approaches]{allocating to alternatives}; \hyperref[meth:land-timber-farmland]{land, timber and farmland}.

\end{qfmethod}
